\section{Organisational Countermeasures}

As established in the previous section, ATP29 combines social engineering with espionage techniques. Their use of relevant material as well as convincing phishing messages suggests that they do not just exploit technical vulnerabilities but also organizational operations \parencite{geers2015cyber, cisa-russia-advisories}. 

Subsequently, these characteristics may have implications for organizational cybersecurity behaviour as employees may be more inclined to trust a message if it is specific to their responsibilities. Furthermore, workplace factors such as productivity and hierarchy may discourage employees from questioning requests. 

Therefore, organizations should not just rely on their employees being capable of identifying suspicious messages but rather develop practices that encourage verification, early reporting, and overall make secure behaviour a part of their daily work. 

Research suggests that employee behaviour is influenced by organizational norms, management, and the working environment, rather than through awareness alone \parencite{Sutton2025}. Similarly, the European Union Agency for Cybersecurity (ENISA) places importance on creating an environment that supports secure behaviour \parencite{europeanSecurity.2018}.  

\subsection{Cybersecurity Culture}

Organizations should prioritize developing a culture where security is understood as an organization-wide responsibility rather than solely the responsibility of the IT or cybersecurity department. This becomes relevant when applied to ATP29 as they use targeted phishing attacks that relate to the professional context of their targets \parencite{geers2015cyber}. Employees may be more inclined to trust such messages when it concerns their responsibilities or appear to come from a trusted source. 

That being said, research on developing cybersecurity culture suggests that management support, security policies and awareness are important elements in establishing a strong culture \parencite{Uchendu2021}. For example, organizations could specify clear expectations that all employees verify unexpected requests for confidential information and refuse to provide credentials, even when these requests appear to come from managers or trusted partners. Additionally, management should encourage employees to take time to verify requests, even if they delay their tasks. 

Such a measure addresses a potential conflict between productivity and security. If employees believe that responding to requests swiftly is more important that to verify them first, they may compromise security for efficiency. On the other hand, if management supports verification and does not criticize employees for doing it, it creates an incentive for employees to act securely.  

\subsection{Targeted Security Training}

Organizations should provide security training based on realistic scenarios rather than generic training that involve unfamiliar senders or obviously suspicious messages that may not sufficiently prepare employees for targeted phishing attacks.  

APT29's use of relevant material suggests that malicious messages may appear legitimate as they target an employee's responsibilities \parencite{geers2015cyber, cisa-russia-advisories}. Furthermore, phishing susceptibility includes additional factors such as individual characteristics, work environment as well as organizational controls rather than just a lack of awareness \parencite{Spicer2026}. 

Therefore, the focus should be on combining general awareness training with specific examples that consider the aforementioned attributes. For example, employees working in IT could practice simulated emails containing requests for credentials of systems access. Additionally, employees should consider whether the request is expected, whether the sender is verified, and whether the requested information should be provided. 

\subsection{Verification and Reporting Procedures}

Organizations should establish straightforward procedures for verifying sensitive requests and reporting suspicious communication. This reduces reliance on the judgement of individual employees and is in line with organizational factors influencing phishing behaviour \parencite{Spicer2026}. On another note, as ATP29 utilizes targeted phishing attacks, malicious messages may appear credible, requiring strong verification capabilities and reporting measures if employees fall for an attack. 

For example, unexpected requests to share confidential information or credentials should be verified through an independent instance. Employees should also be capable of reporting suspicious emails or messages through a designated reporting function or security contact. That being said, organizations should treat reporting as an organizational responsibility rather than a test of whether employees can recognize phishing attacks \parencite{Petri2025}. 

If employees receive convincing messages, they should verify the message and its content before opening any links or sharing sensitive information. If they have already interacted with such a message, they should immediately report the incident. Such a process must be simple, clearly communicated, and most importantly, not result in punishment; otherwise, employees may hesitate, delay seeking help or in the worst-case scenario, not report the incident at all. 

\subsection{Management Support}

Cybersecurity behaviour can be influenced by how management handles employees' decisions and mistakes. For example, if employees expect punishment after an incident, they may attempt to hide the incident which in turn may delay investigation and give attackers more time to exploit compromised systems \parencite{Patterson2024}. 

As stated before, management and organizational support are important aspects of cybersecurity culture while secure behaviour should be integrated into organizational practices \parencite{Sutton2025, europeanSecurity.2018}. 

That being said, organizations should encourage verification, reporting, and distinguish honest mistakes from deliberate violations. For example, if an employee provides sensitive information to a malicious actor, the immediate response should be to report the incident, rather than to assign the blame to the employee, similar to a just culture \parencite{dekker2016just}. This is relevant to ATP29 because convincing communication may lead employees feeling pressured to respond swiftly with an apparently legitimate request.

\subsection{Evaluation}

Organizations should evaluate whether their measures change employee behaviour over time to assess whether their measures result in improved security behaviour. When preparing to defend against parties such as APT29, organizations should evaluate whether employees consistently verify unexpected requests or report suspicious communication.  

On that note, several questions can be considered to evaluate security behaviour: 

\begin{itemize}
    \item \textbf{Reporting rates:} How many employees report suspicious communication? 
    \item \textbf{Verification behaviour:} Do employees verify unexpected requests for sensitive information? 
    \item \textbf{Response time:} How swiftly are suspected incidents reported? 
    \item \textbf{Training effectiveness:} Do employees perform better in later exercises after receiving them at an earlier point of time?
    \item \textbf{Organizational culture:} Do employees believe they can ask for advice and report mistakes without punishment?
\end{itemize}

For example, an organization could conduct an initial exercise, provide training, and repeat a comparable exercise at a later point of time. It could then compare reporting rates, verification behaviour, and the time taken to report suspicious messages. The results would help identify points that require additional attention.~\parencite{europeanSecurity.2018}